\documentclass[12pt]{article}
\usepackage{amsmath, amssymb}
\usepackage{geometry}
\usepackage[dvipsnames,table]{xcolor}
\usepackage{fancyhdr}
\usepackage{tcolorbox}
\usepackage[expansion=false]{microtype}

\tcbuselibrary{skins,breakable}
\geometry{margin=0.9in, top=1in, bottom=1in}

\definecolor{primary}  {RGB}{27,  79,114}
\definecolor{accent}   {RGB}{46,134,193}
\definecolor{lightbg}  {RGB}{234,244,251}
\definecolor{darktext} {RGB}{44, 62, 80}
\definecolor{purpleclr}{RGB}{90, 20,140}
\definecolor{purplebg} {RGB}{240,228,255}

\pagestyle{fancy}\fancyhf{}
\fancyhead[L]{\small\color{primary}\textbf{Transform Theory}\;\textcolor{accent}{|}\;Solved Past Paper 2025}
\fancyhead[R]{\small\color{darktext}BS-IV Mathematics}
\renewcommand{\headrulewidth}{0pt}
\fancyhead[C]{\color{accent}\rule[-1ex]{\textwidth}{1.2pt}}
\fancyfoot[C]{\small\color{darktext}\thepage}

\newtcolorbox{qbox}[1]{enhanced,arc=0pt,boxrule=0pt,
  colback=primary,colframe=primary,left=14pt,right=14pt,top=10pt,bottom=8pt,
  borderline south={4pt}{0pt}{accent},fontupper=\bfseries\color{white},
  title={\small\bfseries\color{accent!80}#1},coltitle=accent!80,
  attach title to upper=\\[2pt]}

\newtcolorbox{ansbox}{enhanced,arc=4pt,boxrule=1.5pt,breakable,
  colback=purplebg!50,colframe=purpleclr,left=10pt,right=10pt,top=8pt,bottom=8pt,
  title={\small\bfseries\color{white}Solution},colbacktitle=purpleclr,coltitle=white,
  attach boxed title to top left={xshift=10pt,yshift=-\tcboxedtitleheight/2},
  boxed title style={arc=3pt,colframe=purpleclr,boxrule=0pt}}

\newcommand{\Lt}{\mathcal{L}}
\newcommand{\Ft}{\mathcal{F}}

\begin{document}

% ── COVER PAGE ─────────────────────────────────────────────
\thispagestyle{empty}
\begin{tcolorbox}[enhanced,arc=0pt,boxrule=0pt,colback=primary,colframe=primary,
  left=20pt,right=20pt,top=28pt,bottom=20pt,borderline south={5pt}{0pt}{accent}]
  \centering
  {\fontsize{26}{32}\bfseries\color{white}Transform Theory\\[8pt]Solved Past Paper}\\[12pt]
  {\large\color{lightbg}Fourier Transforms \textbf{·} Fourier Integrals \textbf{·} Laplace Transforms}
\end{tcolorbox}

\vspace{8pt}
\begin{tcolorbox}[enhanced,arc=0pt,boxrule=0pt,colback=white,colframe=white,
  borderline north={2pt}{0pt}{primary},borderline south={2pt}{0pt}{primary},
  left=0pt,right=0pt,top=0pt,bottom=0pt]
\begin{tabular}{p{0.33\textwidth}|p{0.33\textwidth}|p{0.31\textwidth}}
  \hspace{12pt}\begin{minipage}[t]{0.3\textwidth}\vspace{8pt}
    {\footnotesize\bfseries\color{accent}CLASS}\\[4pt]
    {\bfseries\color{primary}BS-IV \& MSc (Final)}\\[3pt]
    {\footnotesize\color{darktext}Regular Students}
  \vspace{8pt}\end{minipage}&
  \hspace{12pt}\begin{minipage}[t]{0.3\textwidth}\vspace{8pt}
    {\footnotesize\bfseries\color{accent}COMPOSED BY}\\[4pt]
    {\bfseries\color{primary}Nadir Ali Khan}\\[3pt]
    {\footnotesize\color{darktext}BS-IV Mathematics}
  \vspace{8pt}\end{minipage}&
  \hspace{12pt}\begin{minipage}[t]{0.28\textwidth}\vspace{8pt}
    {\footnotesize\bfseries\color{accent}UNIVERSITY}\\[4pt]
    {\bfseries\color{primary}SALU Khairpur}\\[3pt]
    {\footnotesize\color{darktext}Dept.\ of Mathematics}
  \vspace{8pt}\end{minipage}
\end{tabular}
\end{tcolorbox}

\vspace{16pt}
\begin{tcolorbox}[colback=lightbg,colframe=accent,arc=4pt,boxrule=1pt,
  left=12pt,right=12pt,top=10pt,bottom=10pt]
\begin{tabular}{ll}
  \textcolor{accent}{\bfseries Date:}&26 May 2025\\
  \textcolor{accent}{\bfseries Total Marks:}&50\quad(Attempt any \textbf{Four} questions)\\
  \textcolor{accent}{\bfseries Time:}&2 Hours\\
\end{tabular}
\end{tcolorbox}

\newpage

% ══════════════════════════════════════════════════════════
% Q1
% ══════════════════════════════════════════════════════════
\begin{qbox}{Q.No.\ 1 \textemdash\ Solve using Fourier Transform: $\partial^2 u/\partial x^2 = \partial u/\partial t$,\quad $-\infty<x<\infty,\;t>0$}
\end{qbox}

\medskip
\textbf{(i)}\quad $u(x,0)=f(x)=\begin{cases}u_0,&|x|<1\\0,&|x|>1\end{cases}$

\begin{ansbox}
\textbf{Step 1 — Apply Fourier Transform} w.r.t.\ $x$. Let $\hat{u}(\omega,t)=\int_{-\infty}^{\infty}u\,e^{-i\omega x}dx$.
\[
  \Ft\{u_{xx}\}=-\omega^2\hat{u},\qquad \Ft\{u_t\}=\frac{d\hat{u}}{dt}.
\]
The PDE becomes the ODE:
\[
  \frac{d\hat{u}}{dt}+\omega^2\hat{u}=0.
\]
Separate variables: $\dfrac{d\hat{u}}{\hat{u}}=-\omega^2\,dt$. Integrate both sides:
\[
  \ln|\hat{u}|=-\omega^2 t+C \implies \hat{u}(\omega,t)=A(\omega)\,e^{-\omega^2 t}.
\]
Applying $t=0$: $A(\omega)=\hat{u}(\omega,0)$, so $\hat{u}(\omega,t)=\hat{u}(\omega,0)\,e^{-\omega^2 t}$.

\textbf{Step 2 — Find $\hat{u}(\omega,0)$:}
\[
  \hat{u}(\omega,0)=\int_{-1}^{1}u_0\,e^{-i\omega x}dx
  =u_0\!\left[\frac{e^{-i\omega x}}{-i\omega}\right]_{-1}^{1}
  =\frac{u_0}{-i\omega}\left(e^{-i\omega}-e^{i\omega}\right)
  =\frac{2u_0\sin\omega}{\omega}.
\]

\textbf{Step 3 — Solution:}
\[
  \hat{u}(\omega,t)=\frac{2u_0\sin\omega}{\omega}\,e^{-\omega^2 t}.
\]
By inverse Fourier transform:
\[
  \boxed{u(x,t)=\frac{u_0}{\pi}\int_{-\infty}^{\infty}\frac{\sin\omega}{\omega}\,e^{-\omega^2 t+i\omega x}\,d\omega.}
\]
\end{ansbox}

\medskip
\textbf{(ii)}\quad $u(x,0)=e^{-|x|}$

\begin{ansbox}
\textbf{Step 1 — Find $\hat{u}(\omega,0)$:}
\[
  \hat{u}(\omega,0)=\int_{-\infty}^{\infty}e^{-|x|}e^{-i\omega x}dx
  =\int_{-\infty}^{0}e^{x(1-i\omega)}dx+\int_{0}^{\infty}e^{-x(1+i\omega)}dx
  =\frac{1}{1-i\omega}+\frac{1}{1+i\omega}=\frac{2}{1+\omega^2}.
\]

\textbf{Step 2 — Solution:}
\[
  \hat{u}(\omega,t)=\frac{2}{1+\omega^2}\,e^{-\omega^2 t}.
\]
\[
  \boxed{u(x,t)=\frac{1}{\pi}\int_{-\infty}^{\infty}\frac{e^{-\omega^2 t}}{1+\omega^2}\,e^{i\omega x}\,d\omega.}
\]
\end{ansbox}

\vspace{12pt}

% ══════════════════════════════════════════════════════════
% Q2
% ══════════════════════════════════════════════════════════
\begin{qbox}{Q.No.\ 2 \textemdash\ Fourier Cosine and Sine Integral Representations}
\end{qbox}

\medskip
\textit{Recall:}\quad $F_c(\omega)=\int_0^\infty f(x)\cos(\omega x)\,dx$,\quad $F_s(\omega)=\int_0^\infty f(x)\sin(\omega x)\,dx$, and
\[f(x)=\frac{2}{\pi}\int_0^\infty F_c(\omega)\cos(\omega x)\,d\omega=\frac{2}{\pi}\int_0^\infty F_s(\omega)\sin(\omega x)\,d\omega.\]

\medskip\textbf{(i)}\quad $f(x)=xe^{-2x},\;x>0$

\begin{ansbox}
\textbf{Derivation of the formulas} (differentiation under the integral sign):

We know $\displaystyle\int_0^\infty e^{-ax}\cos(bx)\,dx=\frac{a}{a^2+b^2}$. Differentiate both sides w.r.t.\ $a$:
\[
  -\int_0^\infty xe^{-ax}\cos(bx)\,dx=\frac{(a^2+b^2)-2a^2}{(a^2+b^2)^2}=\frac{b^2-a^2}{(a^2+b^2)^2}
  \implies \int_0^\infty xe^{-ax}\cos(bx)\,dx=\frac{a^2-b^2}{(a^2+b^2)^2}.
\]
Similarly from $\displaystyle\int_0^\infty e^{-ax}\sin(bx)\,dx=\frac{b}{a^2+b^2}$, differentiate w.r.t.\ $a$:
\[
  \int_0^\infty xe^{-ax}\sin(bx)\,dx=\frac{2ab}{(a^2+b^2)^2}.
\]
Now apply with $a=2$, $b=\omega$:

\textbf{Cosine transform:}
\[
  F_c(\omega)=\int_0^\infty xe^{-2x}\cos(\omega x)\,dx=\frac{4-\omega^2}{(4+\omega^2)^2}.
\]
\[
  \boxed{xe^{-2x}=\frac{2}{\pi}\int_0^\infty\frac{(4-\omega^2)\cos(\omega x)}{(4+\omega^2)^2}\,d\omega.}
\]

\textbf{Sine transform:}
\[
  F_s(\omega)=\int_0^\infty xe^{-2x}\sin(\omega x)\,dx=\frac{4\omega}{(4+\omega^2)^2}.
\]
\[
  \boxed{xe^{-2x}=\frac{2}{\pi}\int_0^\infty\frac{4\omega\sin(\omega x)}{(4+\omega^2)^2}\,d\omega.}
\]
\end{ansbox}

\medskip\textbf{(ii)}\quad $f(x)=e^{-x}\cos x,\;x>0$

\begin{ansbox}
Use the product-to-sum identity and the result $\int_0^\infty e^{-ax}\cos(bx)\,dx=\dfrac{a}{a^2+b^2}$,\; $\int_0^\infty e^{-ax}\sin(bx)\,dx=\dfrac{b}{a^2+b^2}$.

Write $\cos x\cos(\omega x)=\tfrac{1}{2}[\cos((\omega-1)x)+\cos((\omega+1)x)]$ and $\cos x\sin(\omega x)=\tfrac{1}{2}[\sin((\omega+1)x)+\sin((\omega-1)x)]$.

\textbf{Cosine transform:}
\[
  F_c(\omega)=\frac{1}{2}\!\left[\frac{1}{1+(\omega-1)^2}+\frac{1}{1+(\omega+1)^2}\right].
\]
\[
  \boxed{e^{-x}\cos x=\frac{1}{\pi}\int_0^\infty\!\left[\frac{1}{1+(\omega-1)^2}+\frac{1}{1+(\omega+1)^2}\right]\cos(\omega x)\,d\omega.}
\]

\textbf{Sine transform:}
\[
  F_s(\omega)=\frac{1}{2}\!\left[\frac{\omega+1}{1+(\omega+1)^2}+\frac{\omega-1}{1+(\omega-1)^2}\right].
\]
\[
  \boxed{e^{-x}\cos x=\frac{1}{\pi}\int_0^\infty\!\left[\frac{\omega+1}{1+(\omega+1)^2}+\frac{\omega-1}{1+(\omega-1)^2}\right]\sin(\omega x)\,d\omega.}
\]
\end{ansbox}

\newpage

% ══════════════════════════════════════════════════════════
% Q3
% ══════════════════════════════════════════════════════════
\begin{qbox}{Q.No.\ 3 \textemdash\ Solve $u_{xx}=u_t$ ($x>0,\,t>0$) using Laplace Transform}
\end{qbox}

\medskip\textbf{(i)}\quad $u(0,t)=u_0$,\quad $\lim_{x\to\infty}u(x,t)=u_1$,\quad $u(x,0)=u_1$

\begin{ansbox}
\textbf{Step 1 — Take Laplace} w.r.t.\ $t$. Let $\bar{U}(x,s)=\Lt\{u(x,t)\}$:
\[
  \bar{U}_{xx}-s\bar{U}=-u(x,0)=-u_1 \implies \bar{U}_{xx}-s\bar{U}=-u_1.
\]

\textbf{Step 2 — Solve the ODE in $x$:}

Characteristic equation of the homogeneous part $\bar{U}_{xx}-s\bar{U}=0$:
\[
  r^2-s=0 \implies r=\pm\sqrt{s}.
\]
Homogeneous solution: $\bar{U}_h=Ae^{\sqrt{s}\,x}+Be^{-\sqrt{s}\,x}$.

Particular solution (constant trial $\bar{U}_p=k$): $0-sk=-u_1 \Rightarrow k=u_1/s$. So $\bar{U}_p=u_1/s$.

General: $\bar{U}=Ae^{\sqrt{s}\,x}+Be^{-\sqrt{s}\,x}+u_1/s$.

\textbf{Step 3 — Apply boundary conditions:}

$\lim_{x\to\infty}\bar{U}$ must be finite $\Rightarrow A=0$.

At $x=0$: $\bar{U}(0,s)=\Lt\{u_0\}=u_0/s \Rightarrow B+u_1/s=u_0/s \Rightarrow B=(u_0-u_1)/s$.

\[
  \bar{U}(x,s)=\frac{u_0-u_1}{s}\,e^{-\sqrt{s}\,x}+\frac{u_1}{s}.
\]

\textbf{Step 4 — Inverse Laplace} using $\Lt^{-1}\!\left\{\dfrac{e^{-a\sqrt{s}}}{s}\right\}=\operatorname{erfc}\!\left(\dfrac{a}{2\sqrt{t}}\right)$:
\[
  \boxed{u(x,t)=(u_0-u_1)\operatorname{erfc}\!\left(\frac{x}{2\sqrt{t}}\right)+u_1.}
\]
\textit{Check:} $u(x,0)=(u_0-u_1)\cdot\operatorname{erfc}(\infty)+u_1=u_1$.\;\checkmark\quad $u(0,t)=(u_0-u_1)\cdot1+u_1=u_0$.\;\checkmark
\end{ansbox}

\medskip\textbf{(ii)}\quad $u(0,t)=u_0$,\quad $\displaystyle\lim_{x\to\infty}\frac{u(x,t)}{x}=u_1$,\quad $u(x,0)=u_1 x$

\begin{ansbox}
\textbf{Step 1 — Take Laplace:}
\[
  \bar{U}_{xx}-s\bar{U}=-u(x,0)=-u_1 x.
\]

\textbf{Step 2 — Solve the ODE:}

Characteristic equation: $r^2-s=0\Rightarrow r=\pm\sqrt{s}$, so $\bar{U}_h=Ae^{\sqrt{s}\,x}+Be^{-\sqrt{s}\,x}$.

Particular solution (try $\bar{U}_p=ax+b$):
$-s(ax+b)=-u_1 x \Rightarrow a=u_1/s,\;b=0$. So $\bar{U}_p=u_1 x/s$.

General: $\bar{U}=Ae^{\sqrt{s}\,x}+Be^{-\sqrt{s}\,x}+u_1 x/s$.

\textbf{Step 3 — Apply BCs:}

$\lim_{x\to\infty}\bar{U}/x=u_1/s$ must be finite $\Rightarrow A=0$.

At $x=0$: $\bar{U}(0,s)=u_0/s \Rightarrow B=u_0/s$.
\[
  \bar{U}(x,s)=\frac{u_0}{s}\,e^{-\sqrt{s}\,x}+\frac{u_1 x}{s}.
\]

\textbf{Step 4 — Inverse Laplace:}
\[
  \boxed{u(x,t)=u_0\operatorname{erfc}\!\left(\frac{x}{2\sqrt{t}}\right)+u_1 x.}
\]
\textit{Check:} $u(x,0)=u_0\cdot0+u_1 x=u_1 x$.\;\checkmark\quad $u(0,t)=u_0\cdot1=u_0$.\;\checkmark
\end{ansbox}

\vspace{12pt}

% ══════════════════════════════════════════════════════════
% Q4
% ══════════════════════════════════════════════════════════
\begin{qbox}{Q.No.\ 4 \textemdash\ Inverse Laplace Transform and IVP}
\end{qbox}

\medskip\textbf{(a)}\quad Find $\;\Lt^{-1}\!\left\{\dfrac{s^2+1}{s(s-1)(s+1)(s-2)}\right\}$

\begin{ansbox}
\textbf{Step 1 — Partial fractions:}
\[
  \frac{s^2+1}{s(s-1)(s+1)(s-2)}=\frac{A}{s}+\frac{B}{s-1}+\frac{C}{s+1}+\frac{D}{s-2}.
\]
Cover-up method:
\begin{align*}
  A&=\left.\frac{s^2+1}{(s-1)(s+1)(s-2)}\right|_{s=0}=\frac{1}{(-1)(1)(-2)}=\frac{1}{2},\\
  B&=\left.\frac{s^2+1}{s(s+1)(s-2)}\right|_{s=1}=\frac{2}{(1)(2)(-1)}=-1,\\
  C&=\left.\frac{s^2+1}{s(s-1)(s-2)}\right|_{s=-1}=\frac{2}{(-1)(-2)(-3)}=-\frac{1}{3},\\
  D&=\left.\frac{s^2+1}{s(s-1)(s+1)}\right|_{s=2}=\frac{5}{(2)(1)(3)}=\frac{5}{6}.
\end{align*}

\textbf{Step 2 — Inverse Laplace:}
\[
  \Lt^{-1}\!\left\{\frac{1/2}{s}\right\}-\Lt^{-1}\!\left\{\frac{1}{s-1}\right\}-\Lt^{-1}\!\left\{\frac{1/3}{s+1}\right\}+\Lt^{-1}\!\left\{\frac{5/6}{s-2}\right\}.
\]
\[
  \boxed{\Lt^{-1}\!\left\{\frac{s^2+1}{s(s-1)(s+1)(s-2)}\right\}=\frac{1}{2}-e^t-\frac{1}{3}e^{-t}+\frac{5}{6}e^{2t}.}
\]
\end{ansbox}

\medskip\textbf{(b)}\quad Solve $y''+y=\sqrt{2}\sin(\sqrt{2}\,t)$,\quad $y(0)=10$,\quad $y'(0)=0$

\begin{ansbox}
\textbf{Step 1 — Take Laplace:}
\[
  [s^2Y-10s-0]+Y=\sqrt{2}\cdot\frac{\sqrt{2}}{s^2+2}=\frac{2}{s^2+2}.
\]
\[
  (s^2+1)Y=10s+\frac{2}{s^2+2}\implies Y=\frac{10s}{s^2+1}+\frac{2}{(s^2+1)(s^2+2)}.
\]

\textbf{Step 2 — Partial fractions} for $\dfrac{2}{(s^2+1)(s^2+2)}$. Let $u=s^2$:
\[
  \frac{2}{(u+1)(u+2)}=\frac{A}{u+1}+\frac{B}{u+2}\implies A=2,\;B=-2.
\]
\[
  Y=\frac{10s}{s^2+1}+\frac{2}{s^2+1}-\frac{2}{s^2+2}.
\]

\textbf{Step 3 — Inverse Laplace} (using $\Lt^{-1}\{1/(s^2+a^2)\}=\sin(at)/a$):
\[
  \boxed{y(t)=10\cos t+2\sin t-\sqrt{2}\sin(\sqrt{2}\,t).}
\]
\textit{Check:} $y(0)=10\cdot1+0-0=10$.\;\checkmark\quad $y'(0)=0+2-\sqrt{2}\cdot\sqrt{2}=2-2=0$.\;\checkmark
\end{ansbox}

\vspace{12pt}

% ══════════════════════════════════════════════════════════
% Q5
% ══════════════════════════════════════════════════════════
\begin{qbox}{Q.No.\ 5 \textemdash\ Evaluate using Laplace Transform (do not integrate first)}
\end{qbox}

\medskip\textbf{(a)}\quad $\Lt\{e^{-t}*e^{t}\cos t\}$\quad (convolution)

\begin{ansbox}
By the \textbf{Convolution Theorem}: $\Lt\{f*g\}=F(s)\cdot G(s)$.

Let $f(t)=e^{-t}$ and $g(t)=e^t\cos t$:
\[
  F(s)=\Lt\{e^{-t}\}=\frac{1}{s+1},\qquad G(s)=\Lt\{e^t\cos t\}=\frac{s-1}{(s-1)^2+1}=\frac{s-1}{s^2-2s+2}.
\]
\[
  \boxed{\Lt\{e^{-t}*e^t\cos t\}=\frac{1}{s+1}\cdot\frac{s-1}{s^2-2s+2}=\frac{s-1}{(s+1)(s^2-2s+2)}.}
\]
\end{ansbox}

\medskip\textbf{(b)}\quad $\Lt\!\left\{\displaystyle\int_0^t\sin\tau\cos(t-\tau)\,d\tau\right\}$

\begin{ansbox}
The integral $\displaystyle\int_0^t\sin\tau\cos(t-\tau)\,d\tau$ is the \textbf{convolution} of $\sin t$ and $\cos t$.

By the Convolution Theorem:
\[
  \Lt\{\sin t*\cos t\}=\Lt\{\sin t\}\cdot\Lt\{\cos t\}=\frac{1}{s^2+1}\cdot\frac{s}{s^2+1}.
\]
\[
  \boxed{\Lt\!\left\{\int_0^t\sin\tau\cos(t-\tau)\,d\tau\right\}=\frac{s}{(s^2+1)^2}.}
\]
\end{ansbox}

\end{document}
